\chapter{Control de calidad y muestreo de aceptación}
\label{ch:calidad}

\section{Vigilar un proceso o aceptar un lote}
El control estadístico busca señales de cambio en un proceso; el muestreo de
aceptación decide sobre lotes mediante muestras. Un proceso estable puede
producir defectos, y aceptar un lote no certifica que todas sus unidades sean
conformes. Ambos procedimientos necesitan definiciones y muestreo coherentes
\cite{montgomery2020quality}.

QSB admite cartas c, p y \(\bar X\)-R, Pareto y gráfico de probabilidad
normal; evalúa además planes de aceptación simples y dobles binomiales.
El backend nativo es determinista y el de validación no resuelve. No hay
backend externo disponible para estas familias.

\section{Cartas c y p}
La carta c cuenta defectos en una unidad de inspección de tamaño constante.
Un artículo puede tener más de un defecto. Con hipótesis de conteo Poisson,
\[
 CL=\bar c,\qquad LCL=\max(0,\bar c-3\sqrt{\bar c}),\quad
 UCL=\bar c+3\sqrt{\bar c}.
\]
\path{manual/examples/prob_calidad.json} contiene conteos
\((2,3,4,3,2,4,3,3,2,4)\). Desde la raíz:
\begin{lstlisting}
.build/debug/qsb validate-quality-json manual/examples/prob_calidad.json
.build/debug/qsb solve-quality-json manual/examples/prob_calidad.json
\end{lstlisting}
La línea central es \(3\), límite inferior cero y superior \(8.196152\).
No hay puntos fuera de límites ni señales de las reglas implementadas. Es una
comprobación didáctica de fórmula; diez conteos no bastan para establecer un
proceso industrial como estable. Unidades de inspección de tamaños variables
requieren otro planteamiento; QSB no implementa una carta u en esta ruta.

La carta p utiliza proporción de unidades defectuosas, con tamaño \(n_i\)
para cada subgrupo. Su centro ponderado es
\[
 \bar p=\frac{\sum_i n_ip_i}{\sum_i n_i},\qquad
 CL_i=\bar p,\quad CL_i\pm3\sqrt{\bar p(1-\bar p)/n_i},
\]
recortando límites a \([0,1]\).

En \path{prob_p.json}, tamaños
\((100,100,200)\) y proporciones \((0.02,0.04,0.03)\) producen
\(\bar p=0.03\); los límites superiores son \(0.081176\) para los dos
primeros subgrupos y \(0.066187\) para el tercero. El inferior es cero.
La entrada usa fracciones: \(0.03\) significa \(3\%\), no \(0.03\%\).
Los límites de tres sigmas son una aproximación estadística de la distribución
de los datos, aunque QSB etiquete como \texttt{exact} la aplicación de fórmulas.

\section{Carta de medias y rangos}
Subgrupos racionales agrupan mediciones tomadas bajo condiciones comparables.
Con \(\bar{\bar X}\) media de medias y \(\bar R\) media de rangos,
\[
 LCL_{\bar X},UCL_{\bar X}=\bar{\bar X}\mp A_2\bar R,\qquad
 LCL_R=D_3\bar R,\quad UCL_R=D_4\bar R.
\]
QSB usa constantes estándar para tamaños uniformes entre dos y diez.
Revise primero la variación en R antes de interpretar desplazamientos de medias.
\path{prob_xr.json} contiene subgrupos \((9,10,11)\),
\((10,11,12)\) y \((8,9,10)\). Se obtiene \(\bar{\bar X}=10\),
\(\bar R=2\), límites de medias \([7.954,12.046]\) y de rangos
\([0,5.148]\), sin puntos fuera de límites. Los límites de control describen
variabilidad del proceso; no son tolerancias de diseño. QSB no calcula índices
Cp/Cpk en este flujo.

\section{Pareto y probabilidad normal}
Pareto agrega categorías, las ordena por frecuencia descendente y publica
participaciones y acumulados. \path{prob_pareto.json} tiene filas
\((4,2,1)\) y \((2,1,0)\) para Rayas, Manchas y Roturas. Sus totales son
\((6,3,1)\), total \(10\), participaciones \((0.6,0.3,0.1)\) y acumulados
\((0.6,0.9,1)\). Los campos llamados \texttt{percentage} en JSON almacenan
fracciones, no valores entre cero y cien. Priorizar Rayas por frecuencia no
establece que sea la categoría de mayor costo o riesgo.

El gráfico normal ordena mediciones y asigna posiciones de Blom
\(p_i=(i-3/8)/(n+1/4)\); utiliza \(z_i=\Phi^{-1}(p_i)\), ajustando
medición contra puntuación normal. \path{prob_normal.json} contiene
\((8,9,10,11,12)\); media \(10\), desviación muestral \(1.581139\),
intercepto \(10\), pendiente \(1.742919\) y correlación \(0.997902\).
Se etiqueta \texttt{approximate} por el cuantil normal aproximado. Una
correlación elevada con cinco puntos no demuestra normalidad; no es una prueba
con p-valor. Estos cuatro ejemplos se ejecutan con los mismos comandos de
calidad cambiando el archivo.

\section{Señales e interfaz de calidad}
QSB señala puntos fuera de tres sigmas y reglas fijas Western Electric:
dos de tres más allá de dos sigmas del mismo lado, cuatro de cinco más allá
de una sigma del mismo lado y
ocho del mismo lado del centro. Los puntos fuera de tres sigmas ya señalados
se excluyen del recuento de las reglas zonales de dos de tres y cuatro de cinco.
No hay configuración de reglas ni metadatos
normalizados de causa, acción y comentario. Los índices de señales identifican
observaciones a investigar, no diagnósticos automáticos de causa. Los límites
recortados de cartas de atributos afectan la escala usada por las reglas de
esta implementación: conserve la salida y revise la señal junto al gráfico.

En macOS seleccione la variante en el editor nativo y complete su tabla:
conteos/proporciones, subgrupos, categorías o mediciones. Cambiar variante
conserva tablas no activas durante la sesión, pero el JSON exporta solamente la
activa. Pegue tablas completas para reemplazarlas de forma atómica; datos
incompletos importados permanecen como celdas inválidas. Tras validar y resolver,
las cartas muestran centros, límites y señales; Pareto muestra frecuencia y
acumulado; el gráfico normal muestra puntos y recta. Consulte las tablas y el
JSON para los valores exactos del cálculo y mantenga el orden de adquisición.

\section{Plan simple de aceptación}
Para tamaño de muestra \(n\) y máximo aceptable de defectuosas \(c\), se
acepta si \(X\le c\). Bajo muestreo binomial independiente con proporción
\(p\), la curva característica operativa es
\[
 P_a(p)=\sum_{x=0}^c{n\choose x}p^x(1-p)^{n-x}.
\]
AQL es el nivel de calidad aceptable y RQL el rechazable. El riesgo real del
productor es \(\alpha=1-P_a(AQL)\); el del consumidor es
\(\beta=P_a(RQL)\). Introducir riesgos nominales no modifica un plan ya
suministrado para que los cumpla.

\path{prob_muestreo.json} fija \(n=20,c=1,AQL=0.01,RQL=0.15\),
riesgos nominales \(0.05,0.10\) y lote de \(500\).
\begin{lstlisting}
.build/debug/qsb validate-acceptance-json manual/examples/prob_muestreo.json
.build/debug/qsb solve-acceptance-json manual/examples/prob_muestreo.json
\end{lstlisting}
QSB devuelve \(P_a(AQL)=0.983141\), riesgo del productor \(0.016859\)
y riesgo del consumidor \(0.175558\). El plan supera el riesgo nominal del
consumidor: no debe describirse como un plan que satisface \(\beta=0.10\).
La salida contiene 101 puntos de curva desde \(0\) a \(100\%\) defectuoso,
además de medidas en AQL y RQL. Esta malla no sustituye evaluación continua para
optimizar un riesgo entre puntos.

El tamaño medio de muestra ASN es \(20\). Bajo inspección rectificadora de
lotes rechazados y eliminación de defectuosas inspeccionadas,
\[
 ATI=n+(N-n)(1-P_a),\qquad AOQ=pP_a(N-n)/N.
\]
En AQL resulta ATI \(28.092482\) y AOQ \(0.009438\), es decir
\(0.943815\%\). Si los lotes rechazados se devuelven sin inspección total,
estas medidas no describen el procedimiento real. QSB evalúa las probabilidades
binomiales exactamente dentro del supuesto, pero muestrear sin reemplazo una
fracción grande del lote requiere considerar la distribución hipergeométrica;
el núcleo ofrece esa probabilidad, no un selector en este comando.

\section{Plan doble y diseño}
En un plan doble, se inspeccionan \(n_1\) unidades. Si las defectuosas
\(x_1\le c_1\), se acepta; si \(x_1\ge r_1\), se rechaza; en medio se
inspeccionan \(n_2\) \emph{adicionales}. Se acepta finalmente si
\(x_1+x_2\le c_2\). Confundir \(n_2\) con tamaño acumulado cambia el plan.
Su probabilidad de aceptación es
\[
 P_a=P(X_1\le c_1)+
 \sum_{x=c_1+1}^{r_1-1}P(X_1=x)P(X_2\le c_2-x).
\]
El archivo \path{prob_doble.json} usa \(n_1=10,c_1=0,r_1=3\),
\(n_2=10,c_2=1\) con los mismos niveles y lote anteriores. Cambiando el
archivo en \texttt{solve-acceptance-json}, se obtiene riesgo del productor
\(0.013001\), consumidor \(0.265274\), ASN en AQL \(10.955041\) y en RQL
\(16.233221\). Reduce inspección media a costa de aceptar más lotes malos en
este ejemplo; no es una mejora universal. ATI en AQL es \(17.313145\) y
AOQ \(0.009654\), bajo la convención rectificadora del solver.
Para planes dobles, QSB sustituye \(n\) por ASN en las fórmulas publicadas
de ATI y AOQ. Es una convención agregada: no calcula por separado el tamaño
inspeccionado condicionado a cada ruta de aceptación o rechazo.

La CLI permite diseñar el menor plan \emph{simple} que cumple cotas binomiales:
\begin{lstlisting}
.build/debug/qsb design-acceptance 500 0.01 0.15 0.05 0.10
\end{lstlisting}
El resultado es \(n=25,c=1\), riesgo productor \(0.025759\) y consumidor
\(0.093071\). El orden es lote, AQL, RQL, riesgo productor y riesgo consumidor. El comando
busca un plan, a diferencia de evaluar uno ya fijado. No diseña planes dobles,
no optimiza costo esperado, no implementa errores no nulos de clasificación
y no afirma conformidad con una norma de muestreo específica.

En macOS el formulario distingue plan simple/doble y sus límites de decisión;
conserva los niveles, riesgos y datos económicos compartidos al cambiar tipo.
Valide y resuelva para ver curva OC, marcadores AQL/RQL, riesgos reales frente a
nominales y ASN/ATI/AOQ. Los campos económicos se preservan en el modelo; la
salida del núcleo contiene las medidas de calidad e inspección, no un óptimo
económico. El JSON permite auditar el plan exacto utilizado.

\section{Ejercicios}
\begin{enumerate}
\item Añada un conteo de veinte al ejemplo c y observe cómo cambian centro y
límites al incluirlo en la estimación. Discuta cuándo fijar una fase de referencia.
\item Compare cartas p con tamaño cien y mil manteniendo proporciones.
Explique por qué una misma desviación puede producir señales diferentes.
\item Diseñe un plan simple con las cotas del ejemplo y compare ASN con el
plan \(20,1\). Compruebe riesgos reales, no solo nominales.
\end{enumerate}
